\documentclass[11pt,a4paper]{article}

% shared/macros.tex
% Shared notation for all surface-partition math documents.
% Include with: % shared/macros.tex
% Shared notation for all surface-partition math documents.
% Include with: % shared/macros.tex
% Shared notation for all surface-partition math documents.
% Include with: \input{../shared/macros}

% -------------------------------------------------------------------------
% Packages (include here so each document only needs \input{../shared/macros})
% -------------------------------------------------------------------------
\usepackage{amsmath,amssymb,amsthm}
\usepackage{bm}           % \bm for bold math
\usepackage{mathtools}    % \coloneqq, dcases, etc.
\usepackage{booktabs}     % \toprule etc. in tables
\usepackage{hyperref}
\usepackage{geometry}
\usepackage{microtype}
\usepackage{parskip}
\usepackage{xcolor}
\usepackage{tcolorbox}
\usepackage{enumitem}
\usepackage{float}             % [H] float placement

\geometry{a4paper, margin=2.5cm}

% -------------------------------------------------------------------------
% Theorem-like environments
% -------------------------------------------------------------------------
\theoremstyle{definition}
\newtheorem{defn}{Definition}[section]
\newtheorem{remark}{Remark}[section]

\theoremstyle{plain}
\newtheorem{prop}{Proposition}[section]

% -------------------------------------------------------------------------
% Mesh and geometry
% -------------------------------------------------------------------------
\newcommand{\V}{\mathbf{V}}            % vertex array V[i] in R^dim
\newcommand{\F}{\mathbf{F}}            % face (triangle) array
\newcommand{\vone}[1]{v_{1,#1}}       % first vertex index of VP i
\newcommand{\vtwo}[1]{v_{2,#1}}       % second vertex index of VP i

% -------------------------------------------------------------------------
% Variable points
% -------------------------------------------------------------------------
\newcommand{\lam}{\lambda}             % scalar VP parameter
\newcommand{\Lam}{\boldsymbol{\lambda}}% full parameter vector
\newcommand{\vp}[1]{\mathbf{p}_{#1}}  % 3D position of VP i
\newcommand{\ed}[1]{\mathbf{d}_{#1}}  % edge direction d_i = V[v1_i] - V[v2_i]

% -------------------------------------------------------------------------
% Steiner / triple-point
% -------------------------------------------------------------------------
\newcommand{\Sv}{\mathbf{S}}           % Steiner (Fermat-Torricelli) point
\newcommand{\Ki}{K_{i}}               % projection matrix (I-n_i n_i^T)/r_i
\newcommand{\Mmat}{M}                 % sum of K_i matrices at a triple point

% -------------------------------------------------------------------------
% Normals and unit vectors
% -------------------------------------------------------------------------
\newcommand{\nhat}{\hat{\mathbf{n}}}  % unit normal to a triangle
\newcommand{\Dhat}{\hat{\boldsymbol{\Delta}}}  % unit segment direction

% -------------------------------------------------------------------------
% Operations
% -------------------------------------------------------------------------
\newcommand{\norm}[1]{\left\|#1\right\|}
\newcommand{\abs}[1]{\left|#1\right|}
\newcommand{\ip}[2]{\langle #1,\, #2 \rangle}

% Projection perpendicular to a unit vector \uhat
\newcommand{\Proj}[1]{\bigl(\mathbf{I} - #1 #1^\top\bigr)}

% argmax operator (winner-take-all assignment in Phase 1 documents)
\DeclareMathOperator*{\argmax}{arg\,max}
\DeclareMathOperator*{\argmin}{arg\,min}

% -------------------------------------------------------------------------
% Calculus shorthands
% -------------------------------------------------------------------------
\newcommand{\pd}[2]{\dfrac{\partial #1}{\partial #2}}
\newcommand{\pdd}[3]{\dfrac{\partial^2 #1}{\partial #2\,\partial #3}}
\newcommand{\grad}{\nabla}
\newcommand{\Hess}[1]{\nabla^2 #1}

% -------------------------------------------------------------------------
% Optimization problem objects
% -------------------------------------------------------------------------
\newcommand{\Preg}{P_{\mathrm{reg}}}  % regular perimeter
\newcommand{\Pst}{P_{\mathrm{st}}}   % Steiner perimeter
\newcommand{\Areg}{A^{\mathrm{reg}}} % regular cell area
\newcommand{\Ast}{A^{\mathrm{st}}}   % Steiner area contribution
\newcommand{\Atgt}{\bar{A}}           % target area per cell
\newcommand{\Lagr}{\mathcal{L}}       % Lagrangian
\newcommand{\objfac}{\sigma}          % IPOPT objective scaling factor

% -------------------------------------------------------------------------
% Code cross-references (for "In the code..." callouts)
% -------------------------------------------------------------------------
\newcommand{\code}[1]{\texttt{#1}}
\newcommand{\codefile}[1]{\texttt{\small #1}}

% -------------------------------------------------------------------------
% Threshold dynamics / approach B (10-mbo-auction-dynamics)
% -------------------------------------------------------------------------
\newcommand{\Dlump}{D}                 % lumped mass matrix diag(v)
\newcommand{\Atau}{A_\tau}             % discrete diffusion operator (M + tau K)^{-1} M
\newcommand{\Etau}{E_\tau}             % heat-content (threshold-dynamics) energy
\newcommand{\Rcell}{R_{\mathrm{cell}}} % radius of the equal-area disc sqrt(Abar/pi)
\newcommand{\hmean}{h_{\mathrm{mean}}} % mean edge length
\newcommand{\hmax}{h_{\mathrm{max}}}   % max edge length
\newcommand{\Gt}{G_t}                  % heat kernel at time t
\newcommand{\Ktau}{K^{\mathrm{res}}_\tau} % resolvent (backward-Euler) kernel


% -------------------------------------------------------------------------
% Packages (include here so each document only needs % shared/macros.tex
% Shared notation for all surface-partition math documents.
% Include with: \input{../shared/macros}

% -------------------------------------------------------------------------
% Packages (include here so each document only needs \input{../shared/macros})
% -------------------------------------------------------------------------
\usepackage{amsmath,amssymb,amsthm}
\usepackage{bm}           % \bm for bold math
\usepackage{mathtools}    % \coloneqq, dcases, etc.
\usepackage{booktabs}     % \toprule etc. in tables
\usepackage{hyperref}
\usepackage{geometry}
\usepackage{microtype}
\usepackage{parskip}
\usepackage{xcolor}
\usepackage{tcolorbox}
\usepackage{enumitem}
\usepackage{float}             % [H] float placement

\geometry{a4paper, margin=2.5cm}

% -------------------------------------------------------------------------
% Theorem-like environments
% -------------------------------------------------------------------------
\theoremstyle{definition}
\newtheorem{defn}{Definition}[section]
\newtheorem{remark}{Remark}[section]

\theoremstyle{plain}
\newtheorem{prop}{Proposition}[section]

% -------------------------------------------------------------------------
% Mesh and geometry
% -------------------------------------------------------------------------
\newcommand{\V}{\mathbf{V}}            % vertex array V[i] in R^dim
\newcommand{\F}{\mathbf{F}}            % face (triangle) array
\newcommand{\vone}[1]{v_{1,#1}}       % first vertex index of VP i
\newcommand{\vtwo}[1]{v_{2,#1}}       % second vertex index of VP i

% -------------------------------------------------------------------------
% Variable points
% -------------------------------------------------------------------------
\newcommand{\lam}{\lambda}             % scalar VP parameter
\newcommand{\Lam}{\boldsymbol{\lambda}}% full parameter vector
\newcommand{\vp}[1]{\mathbf{p}_{#1}}  % 3D position of VP i
\newcommand{\ed}[1]{\mathbf{d}_{#1}}  % edge direction d_i = V[v1_i] - V[v2_i]

% -------------------------------------------------------------------------
% Steiner / triple-point
% -------------------------------------------------------------------------
\newcommand{\Sv}{\mathbf{S}}           % Steiner (Fermat-Torricelli) point
\newcommand{\Ki}{K_{i}}               % projection matrix (I-n_i n_i^T)/r_i
\newcommand{\Mmat}{M}                 % sum of K_i matrices at a triple point

% -------------------------------------------------------------------------
% Normals and unit vectors
% -------------------------------------------------------------------------
\newcommand{\nhat}{\hat{\mathbf{n}}}  % unit normal to a triangle
\newcommand{\Dhat}{\hat{\boldsymbol{\Delta}}}  % unit segment direction

% -------------------------------------------------------------------------
% Operations
% -------------------------------------------------------------------------
\newcommand{\norm}[1]{\left\|#1\right\|}
\newcommand{\abs}[1]{\left|#1\right|}
\newcommand{\ip}[2]{\langle #1,\, #2 \rangle}

% Projection perpendicular to a unit vector \uhat
\newcommand{\Proj}[1]{\bigl(\mathbf{I} - #1 #1^\top\bigr)}

% argmax operator (winner-take-all assignment in Phase 1 documents)
\DeclareMathOperator*{\argmax}{arg\,max}
\DeclareMathOperator*{\argmin}{arg\,min}

% -------------------------------------------------------------------------
% Calculus shorthands
% -------------------------------------------------------------------------
\newcommand{\pd}[2]{\dfrac{\partial #1}{\partial #2}}
\newcommand{\pdd}[3]{\dfrac{\partial^2 #1}{\partial #2\,\partial #3}}
\newcommand{\grad}{\nabla}
\newcommand{\Hess}[1]{\nabla^2 #1}

% -------------------------------------------------------------------------
% Optimization problem objects
% -------------------------------------------------------------------------
\newcommand{\Preg}{P_{\mathrm{reg}}}  % regular perimeter
\newcommand{\Pst}{P_{\mathrm{st}}}   % Steiner perimeter
\newcommand{\Areg}{A^{\mathrm{reg}}} % regular cell area
\newcommand{\Ast}{A^{\mathrm{st}}}   % Steiner area contribution
\newcommand{\Atgt}{\bar{A}}           % target area per cell
\newcommand{\Lagr}{\mathcal{L}}       % Lagrangian
\newcommand{\objfac}{\sigma}          % IPOPT objective scaling factor

% -------------------------------------------------------------------------
% Code cross-references (for "In the code..." callouts)
% -------------------------------------------------------------------------
\newcommand{\code}[1]{\texttt{#1}}
\newcommand{\codefile}[1]{\texttt{\small #1}}

% -------------------------------------------------------------------------
% Threshold dynamics / approach B (10-mbo-auction-dynamics)
% -------------------------------------------------------------------------
\newcommand{\Dlump}{D}                 % lumped mass matrix diag(v)
\newcommand{\Atau}{A_\tau}             % discrete diffusion operator (M + tau K)^{-1} M
\newcommand{\Etau}{E_\tau}             % heat-content (threshold-dynamics) energy
\newcommand{\Rcell}{R_{\mathrm{cell}}} % radius of the equal-area disc sqrt(Abar/pi)
\newcommand{\hmean}{h_{\mathrm{mean}}} % mean edge length
\newcommand{\hmax}{h_{\mathrm{max}}}   % max edge length
\newcommand{\Gt}{G_t}                  % heat kernel at time t
\newcommand{\Ktau}{K^{\mathrm{res}}_\tau} % resolvent (backward-Euler) kernel
)
% -------------------------------------------------------------------------
\usepackage{amsmath,amssymb,amsthm}
\usepackage{bm}           % \bm for bold math
\usepackage{mathtools}    % \coloneqq, dcases, etc.
\usepackage{booktabs}     % \toprule etc. in tables
\usepackage{hyperref}
\usepackage{geometry}
\usepackage{microtype}
\usepackage{parskip}
\usepackage{xcolor}
\usepackage{tcolorbox}
\usepackage{enumitem}
\usepackage{float}             % [H] float placement

\geometry{a4paper, margin=2.5cm}

% -------------------------------------------------------------------------
% Theorem-like environments
% -------------------------------------------------------------------------
\theoremstyle{definition}
\newtheorem{defn}{Definition}[section]
\newtheorem{remark}{Remark}[section]

\theoremstyle{plain}
\newtheorem{prop}{Proposition}[section]

% -------------------------------------------------------------------------
% Mesh and geometry
% -------------------------------------------------------------------------
\newcommand{\V}{\mathbf{V}}            % vertex array V[i] in R^dim
\newcommand{\F}{\mathbf{F}}            % face (triangle) array
\newcommand{\vone}[1]{v_{1,#1}}       % first vertex index of VP i
\newcommand{\vtwo}[1]{v_{2,#1}}       % second vertex index of VP i

% -------------------------------------------------------------------------
% Variable points
% -------------------------------------------------------------------------
\newcommand{\lam}{\lambda}             % scalar VP parameter
\newcommand{\Lam}{\boldsymbol{\lambda}}% full parameter vector
\newcommand{\vp}[1]{\mathbf{p}_{#1}}  % 3D position of VP i
\newcommand{\ed}[1]{\mathbf{d}_{#1}}  % edge direction d_i = V[v1_i] - V[v2_i]

% -------------------------------------------------------------------------
% Steiner / triple-point
% -------------------------------------------------------------------------
\newcommand{\Sv}{\mathbf{S}}           % Steiner (Fermat-Torricelli) point
\newcommand{\Ki}{K_{i}}               % projection matrix (I-n_i n_i^T)/r_i
\newcommand{\Mmat}{M}                 % sum of K_i matrices at a triple point

% -------------------------------------------------------------------------
% Normals and unit vectors
% -------------------------------------------------------------------------
\newcommand{\nhat}{\hat{\mathbf{n}}}  % unit normal to a triangle
\newcommand{\Dhat}{\hat{\boldsymbol{\Delta}}}  % unit segment direction

% -------------------------------------------------------------------------
% Operations
% -------------------------------------------------------------------------
\newcommand{\norm}[1]{\left\|#1\right\|}
\newcommand{\abs}[1]{\left|#1\right|}
\newcommand{\ip}[2]{\langle #1,\, #2 \rangle}

% Projection perpendicular to a unit vector \uhat
\newcommand{\Proj}[1]{\bigl(\mathbf{I} - #1 #1^\top\bigr)}

% argmax operator (winner-take-all assignment in Phase 1 documents)
\DeclareMathOperator*{\argmax}{arg\,max}
\DeclareMathOperator*{\argmin}{arg\,min}

% -------------------------------------------------------------------------
% Calculus shorthands
% -------------------------------------------------------------------------
\newcommand{\pd}[2]{\dfrac{\partial #1}{\partial #2}}
\newcommand{\pdd}[3]{\dfrac{\partial^2 #1}{\partial #2\,\partial #3}}
\newcommand{\grad}{\nabla}
\newcommand{\Hess}[1]{\nabla^2 #1}

% -------------------------------------------------------------------------
% Optimization problem objects
% -------------------------------------------------------------------------
\newcommand{\Preg}{P_{\mathrm{reg}}}  % regular perimeter
\newcommand{\Pst}{P_{\mathrm{st}}}   % Steiner perimeter
\newcommand{\Areg}{A^{\mathrm{reg}}} % regular cell area
\newcommand{\Ast}{A^{\mathrm{st}}}   % Steiner area contribution
\newcommand{\Atgt}{\bar{A}}           % target area per cell
\newcommand{\Lagr}{\mathcal{L}}       % Lagrangian
\newcommand{\objfac}{\sigma}          % IPOPT objective scaling factor

% -------------------------------------------------------------------------
% Code cross-references (for "In the code..." callouts)
% -------------------------------------------------------------------------
\newcommand{\code}[1]{\texttt{#1}}
\newcommand{\codefile}[1]{\texttt{\small #1}}

% -------------------------------------------------------------------------
% Threshold dynamics / approach B (10-mbo-auction-dynamics)
% -------------------------------------------------------------------------
\newcommand{\Dlump}{D}                 % lumped mass matrix diag(v)
\newcommand{\Atau}{A_\tau}             % discrete diffusion operator (M + tau K)^{-1} M
\newcommand{\Etau}{E_\tau}             % heat-content (threshold-dynamics) energy
\newcommand{\Rcell}{R_{\mathrm{cell}}} % radius of the equal-area disc sqrt(Abar/pi)
\newcommand{\hmean}{h_{\mathrm{mean}}} % mean edge length
\newcommand{\hmax}{h_{\mathrm{max}}}   % max edge length
\newcommand{\Gt}{G_t}                  % heat kernel at time t
\newcommand{\Ktau}{K^{\mathrm{res}}_\tau} % resolvent (backward-Euler) kernel


% -------------------------------------------------------------------------
% Packages (include here so each document only needs % shared/macros.tex
% Shared notation for all surface-partition math documents.
% Include with: % shared/macros.tex
% Shared notation for all surface-partition math documents.
% Include with: \input{../shared/macros}

% -------------------------------------------------------------------------
% Packages (include here so each document only needs \input{../shared/macros})
% -------------------------------------------------------------------------
\usepackage{amsmath,amssymb,amsthm}
\usepackage{bm}           % \bm for bold math
\usepackage{mathtools}    % \coloneqq, dcases, etc.
\usepackage{booktabs}     % \toprule etc. in tables
\usepackage{hyperref}
\usepackage{geometry}
\usepackage{microtype}
\usepackage{parskip}
\usepackage{xcolor}
\usepackage{tcolorbox}
\usepackage{enumitem}
\usepackage{float}             % [H] float placement

\geometry{a4paper, margin=2.5cm}

% -------------------------------------------------------------------------
% Theorem-like environments
% -------------------------------------------------------------------------
\theoremstyle{definition}
\newtheorem{defn}{Definition}[section]
\newtheorem{remark}{Remark}[section]

\theoremstyle{plain}
\newtheorem{prop}{Proposition}[section]

% -------------------------------------------------------------------------
% Mesh and geometry
% -------------------------------------------------------------------------
\newcommand{\V}{\mathbf{V}}            % vertex array V[i] in R^dim
\newcommand{\F}{\mathbf{F}}            % face (triangle) array
\newcommand{\vone}[1]{v_{1,#1}}       % first vertex index of VP i
\newcommand{\vtwo}[1]{v_{2,#1}}       % second vertex index of VP i

% -------------------------------------------------------------------------
% Variable points
% -------------------------------------------------------------------------
\newcommand{\lam}{\lambda}             % scalar VP parameter
\newcommand{\Lam}{\boldsymbol{\lambda}}% full parameter vector
\newcommand{\vp}[1]{\mathbf{p}_{#1}}  % 3D position of VP i
\newcommand{\ed}[1]{\mathbf{d}_{#1}}  % edge direction d_i = V[v1_i] - V[v2_i]

% -------------------------------------------------------------------------
% Steiner / triple-point
% -------------------------------------------------------------------------
\newcommand{\Sv}{\mathbf{S}}           % Steiner (Fermat-Torricelli) point
\newcommand{\Ki}{K_{i}}               % projection matrix (I-n_i n_i^T)/r_i
\newcommand{\Mmat}{M}                 % sum of K_i matrices at a triple point

% -------------------------------------------------------------------------
% Normals and unit vectors
% -------------------------------------------------------------------------
\newcommand{\nhat}{\hat{\mathbf{n}}}  % unit normal to a triangle
\newcommand{\Dhat}{\hat{\boldsymbol{\Delta}}}  % unit segment direction

% -------------------------------------------------------------------------
% Operations
% -------------------------------------------------------------------------
\newcommand{\norm}[1]{\left\|#1\right\|}
\newcommand{\abs}[1]{\left|#1\right|}
\newcommand{\ip}[2]{\langle #1,\, #2 \rangle}

% Projection perpendicular to a unit vector \uhat
\newcommand{\Proj}[1]{\bigl(\mathbf{I} - #1 #1^\top\bigr)}

% argmax operator (winner-take-all assignment in Phase 1 documents)
\DeclareMathOperator*{\argmax}{arg\,max}
\DeclareMathOperator*{\argmin}{arg\,min}

% -------------------------------------------------------------------------
% Calculus shorthands
% -------------------------------------------------------------------------
\newcommand{\pd}[2]{\dfrac{\partial #1}{\partial #2}}
\newcommand{\pdd}[3]{\dfrac{\partial^2 #1}{\partial #2\,\partial #3}}
\newcommand{\grad}{\nabla}
\newcommand{\Hess}[1]{\nabla^2 #1}

% -------------------------------------------------------------------------
% Optimization problem objects
% -------------------------------------------------------------------------
\newcommand{\Preg}{P_{\mathrm{reg}}}  % regular perimeter
\newcommand{\Pst}{P_{\mathrm{st}}}   % Steiner perimeter
\newcommand{\Areg}{A^{\mathrm{reg}}} % regular cell area
\newcommand{\Ast}{A^{\mathrm{st}}}   % Steiner area contribution
\newcommand{\Atgt}{\bar{A}}           % target area per cell
\newcommand{\Lagr}{\mathcal{L}}       % Lagrangian
\newcommand{\objfac}{\sigma}          % IPOPT objective scaling factor

% -------------------------------------------------------------------------
% Code cross-references (for "In the code..." callouts)
% -------------------------------------------------------------------------
\newcommand{\code}[1]{\texttt{#1}}
\newcommand{\codefile}[1]{\texttt{\small #1}}

% -------------------------------------------------------------------------
% Threshold dynamics / approach B (10-mbo-auction-dynamics)
% -------------------------------------------------------------------------
\newcommand{\Dlump}{D}                 % lumped mass matrix diag(v)
\newcommand{\Atau}{A_\tau}             % discrete diffusion operator (M + tau K)^{-1} M
\newcommand{\Etau}{E_\tau}             % heat-content (threshold-dynamics) energy
\newcommand{\Rcell}{R_{\mathrm{cell}}} % radius of the equal-area disc sqrt(Abar/pi)
\newcommand{\hmean}{h_{\mathrm{mean}}} % mean edge length
\newcommand{\hmax}{h_{\mathrm{max}}}   % max edge length
\newcommand{\Gt}{G_t}                  % heat kernel at time t
\newcommand{\Ktau}{K^{\mathrm{res}}_\tau} % resolvent (backward-Euler) kernel


% -------------------------------------------------------------------------
% Packages (include here so each document only needs % shared/macros.tex
% Shared notation for all surface-partition math documents.
% Include with: \input{../shared/macros}

% -------------------------------------------------------------------------
% Packages (include here so each document only needs \input{../shared/macros})
% -------------------------------------------------------------------------
\usepackage{amsmath,amssymb,amsthm}
\usepackage{bm}           % \bm for bold math
\usepackage{mathtools}    % \coloneqq, dcases, etc.
\usepackage{booktabs}     % \toprule etc. in tables
\usepackage{hyperref}
\usepackage{geometry}
\usepackage{microtype}
\usepackage{parskip}
\usepackage{xcolor}
\usepackage{tcolorbox}
\usepackage{enumitem}
\usepackage{float}             % [H] float placement

\geometry{a4paper, margin=2.5cm}

% -------------------------------------------------------------------------
% Theorem-like environments
% -------------------------------------------------------------------------
\theoremstyle{definition}
\newtheorem{defn}{Definition}[section]
\newtheorem{remark}{Remark}[section]

\theoremstyle{plain}
\newtheorem{prop}{Proposition}[section]

% -------------------------------------------------------------------------
% Mesh and geometry
% -------------------------------------------------------------------------
\newcommand{\V}{\mathbf{V}}            % vertex array V[i] in R^dim
\newcommand{\F}{\mathbf{F}}            % face (triangle) array
\newcommand{\vone}[1]{v_{1,#1}}       % first vertex index of VP i
\newcommand{\vtwo}[1]{v_{2,#1}}       % second vertex index of VP i

% -------------------------------------------------------------------------
% Variable points
% -------------------------------------------------------------------------
\newcommand{\lam}{\lambda}             % scalar VP parameter
\newcommand{\Lam}{\boldsymbol{\lambda}}% full parameter vector
\newcommand{\vp}[1]{\mathbf{p}_{#1}}  % 3D position of VP i
\newcommand{\ed}[1]{\mathbf{d}_{#1}}  % edge direction d_i = V[v1_i] - V[v2_i]

% -------------------------------------------------------------------------
% Steiner / triple-point
% -------------------------------------------------------------------------
\newcommand{\Sv}{\mathbf{S}}           % Steiner (Fermat-Torricelli) point
\newcommand{\Ki}{K_{i}}               % projection matrix (I-n_i n_i^T)/r_i
\newcommand{\Mmat}{M}                 % sum of K_i matrices at a triple point

% -------------------------------------------------------------------------
% Normals and unit vectors
% -------------------------------------------------------------------------
\newcommand{\nhat}{\hat{\mathbf{n}}}  % unit normal to a triangle
\newcommand{\Dhat}{\hat{\boldsymbol{\Delta}}}  % unit segment direction

% -------------------------------------------------------------------------
% Operations
% -------------------------------------------------------------------------
\newcommand{\norm}[1]{\left\|#1\right\|}
\newcommand{\abs}[1]{\left|#1\right|}
\newcommand{\ip}[2]{\langle #1,\, #2 \rangle}

% Projection perpendicular to a unit vector \uhat
\newcommand{\Proj}[1]{\bigl(\mathbf{I} - #1 #1^\top\bigr)}

% argmax operator (winner-take-all assignment in Phase 1 documents)
\DeclareMathOperator*{\argmax}{arg\,max}
\DeclareMathOperator*{\argmin}{arg\,min}

% -------------------------------------------------------------------------
% Calculus shorthands
% -------------------------------------------------------------------------
\newcommand{\pd}[2]{\dfrac{\partial #1}{\partial #2}}
\newcommand{\pdd}[3]{\dfrac{\partial^2 #1}{\partial #2\,\partial #3}}
\newcommand{\grad}{\nabla}
\newcommand{\Hess}[1]{\nabla^2 #1}

% -------------------------------------------------------------------------
% Optimization problem objects
% -------------------------------------------------------------------------
\newcommand{\Preg}{P_{\mathrm{reg}}}  % regular perimeter
\newcommand{\Pst}{P_{\mathrm{st}}}   % Steiner perimeter
\newcommand{\Areg}{A^{\mathrm{reg}}} % regular cell area
\newcommand{\Ast}{A^{\mathrm{st}}}   % Steiner area contribution
\newcommand{\Atgt}{\bar{A}}           % target area per cell
\newcommand{\Lagr}{\mathcal{L}}       % Lagrangian
\newcommand{\objfac}{\sigma}          % IPOPT objective scaling factor

% -------------------------------------------------------------------------
% Code cross-references (for "In the code..." callouts)
% -------------------------------------------------------------------------
\newcommand{\code}[1]{\texttt{#1}}
\newcommand{\codefile}[1]{\texttt{\small #1}}

% -------------------------------------------------------------------------
% Threshold dynamics / approach B (10-mbo-auction-dynamics)
% -------------------------------------------------------------------------
\newcommand{\Dlump}{D}                 % lumped mass matrix diag(v)
\newcommand{\Atau}{A_\tau}             % discrete diffusion operator (M + tau K)^{-1} M
\newcommand{\Etau}{E_\tau}             % heat-content (threshold-dynamics) energy
\newcommand{\Rcell}{R_{\mathrm{cell}}} % radius of the equal-area disc sqrt(Abar/pi)
\newcommand{\hmean}{h_{\mathrm{mean}}} % mean edge length
\newcommand{\hmax}{h_{\mathrm{max}}}   % max edge length
\newcommand{\Gt}{G_t}                  % heat kernel at time t
\newcommand{\Ktau}{K^{\mathrm{res}}_\tau} % resolvent (backward-Euler) kernel
)
% -------------------------------------------------------------------------
\usepackage{amsmath,amssymb,amsthm}
\usepackage{bm}           % \bm for bold math
\usepackage{mathtools}    % \coloneqq, dcases, etc.
\usepackage{booktabs}     % \toprule etc. in tables
\usepackage{hyperref}
\usepackage{geometry}
\usepackage{microtype}
\usepackage{parskip}
\usepackage{xcolor}
\usepackage{tcolorbox}
\usepackage{enumitem}
\usepackage{float}             % [H] float placement

\geometry{a4paper, margin=2.5cm}

% -------------------------------------------------------------------------
% Theorem-like environments
% -------------------------------------------------------------------------
\theoremstyle{definition}
\newtheorem{defn}{Definition}[section]
\newtheorem{remark}{Remark}[section]

\theoremstyle{plain}
\newtheorem{prop}{Proposition}[section]

% -------------------------------------------------------------------------
% Mesh and geometry
% -------------------------------------------------------------------------
\newcommand{\V}{\mathbf{V}}            % vertex array V[i] in R^dim
\newcommand{\F}{\mathbf{F}}            % face (triangle) array
\newcommand{\vone}[1]{v_{1,#1}}       % first vertex index of VP i
\newcommand{\vtwo}[1]{v_{2,#1}}       % second vertex index of VP i

% -------------------------------------------------------------------------
% Variable points
% -------------------------------------------------------------------------
\newcommand{\lam}{\lambda}             % scalar VP parameter
\newcommand{\Lam}{\boldsymbol{\lambda}}% full parameter vector
\newcommand{\vp}[1]{\mathbf{p}_{#1}}  % 3D position of VP i
\newcommand{\ed}[1]{\mathbf{d}_{#1}}  % edge direction d_i = V[v1_i] - V[v2_i]

% -------------------------------------------------------------------------
% Steiner / triple-point
% -------------------------------------------------------------------------
\newcommand{\Sv}{\mathbf{S}}           % Steiner (Fermat-Torricelli) point
\newcommand{\Ki}{K_{i}}               % projection matrix (I-n_i n_i^T)/r_i
\newcommand{\Mmat}{M}                 % sum of K_i matrices at a triple point

% -------------------------------------------------------------------------
% Normals and unit vectors
% -------------------------------------------------------------------------
\newcommand{\nhat}{\hat{\mathbf{n}}}  % unit normal to a triangle
\newcommand{\Dhat}{\hat{\boldsymbol{\Delta}}}  % unit segment direction

% -------------------------------------------------------------------------
% Operations
% -------------------------------------------------------------------------
\newcommand{\norm}[1]{\left\|#1\right\|}
\newcommand{\abs}[1]{\left|#1\right|}
\newcommand{\ip}[2]{\langle #1,\, #2 \rangle}

% Projection perpendicular to a unit vector \uhat
\newcommand{\Proj}[1]{\bigl(\mathbf{I} - #1 #1^\top\bigr)}

% argmax operator (winner-take-all assignment in Phase 1 documents)
\DeclareMathOperator*{\argmax}{arg\,max}
\DeclareMathOperator*{\argmin}{arg\,min}

% -------------------------------------------------------------------------
% Calculus shorthands
% -------------------------------------------------------------------------
\newcommand{\pd}[2]{\dfrac{\partial #1}{\partial #2}}
\newcommand{\pdd}[3]{\dfrac{\partial^2 #1}{\partial #2\,\partial #3}}
\newcommand{\grad}{\nabla}
\newcommand{\Hess}[1]{\nabla^2 #1}

% -------------------------------------------------------------------------
% Optimization problem objects
% -------------------------------------------------------------------------
\newcommand{\Preg}{P_{\mathrm{reg}}}  % regular perimeter
\newcommand{\Pst}{P_{\mathrm{st}}}   % Steiner perimeter
\newcommand{\Areg}{A^{\mathrm{reg}}} % regular cell area
\newcommand{\Ast}{A^{\mathrm{st}}}   % Steiner area contribution
\newcommand{\Atgt}{\bar{A}}           % target area per cell
\newcommand{\Lagr}{\mathcal{L}}       % Lagrangian
\newcommand{\objfac}{\sigma}          % IPOPT objective scaling factor

% -------------------------------------------------------------------------
% Code cross-references (for "In the code..." callouts)
% -------------------------------------------------------------------------
\newcommand{\code}[1]{\texttt{#1}}
\newcommand{\codefile}[1]{\texttt{\small #1}}

% -------------------------------------------------------------------------
% Threshold dynamics / approach B (10-mbo-auction-dynamics)
% -------------------------------------------------------------------------
\newcommand{\Dlump}{D}                 % lumped mass matrix diag(v)
\newcommand{\Atau}{A_\tau}             % discrete diffusion operator (M + tau K)^{-1} M
\newcommand{\Etau}{E_\tau}             % heat-content (threshold-dynamics) energy
\newcommand{\Rcell}{R_{\mathrm{cell}}} % radius of the equal-area disc sqrt(Abar/pi)
\newcommand{\hmean}{h_{\mathrm{mean}}} % mean edge length
\newcommand{\hmax}{h_{\mathrm{max}}}   % max edge length
\newcommand{\Gt}{G_t}                  % heat kernel at time t
\newcommand{\Ktau}{K^{\mathrm{res}}_\tau} % resolvent (backward-Euler) kernel
)
% -------------------------------------------------------------------------
\usepackage{amsmath,amssymb,amsthm}
\usepackage{bm}           % \bm for bold math
\usepackage{mathtools}    % \coloneqq, dcases, etc.
\usepackage{booktabs}     % \toprule etc. in tables
\usepackage{hyperref}
\usepackage{geometry}
\usepackage{microtype}
\usepackage{parskip}
\usepackage{xcolor}
\usepackage{tcolorbox}
\usepackage{enumitem}
\usepackage{float}             % [H] float placement

\geometry{a4paper, margin=2.5cm}

% -------------------------------------------------------------------------
% Theorem-like environments
% -------------------------------------------------------------------------
\theoremstyle{definition}
\newtheorem{defn}{Definition}[section]
\newtheorem{remark}{Remark}[section]

\theoremstyle{plain}
\newtheorem{prop}{Proposition}[section]

% -------------------------------------------------------------------------
% Mesh and geometry
% -------------------------------------------------------------------------
\newcommand{\V}{\mathbf{V}}            % vertex array V[i] in R^dim
\newcommand{\F}{\mathbf{F}}            % face (triangle) array
\newcommand{\vone}[1]{v_{1,#1}}       % first vertex index of VP i
\newcommand{\vtwo}[1]{v_{2,#1}}       % second vertex index of VP i

% -------------------------------------------------------------------------
% Variable points
% -------------------------------------------------------------------------
\newcommand{\lam}{\lambda}             % scalar VP parameter
\newcommand{\Lam}{\boldsymbol{\lambda}}% full parameter vector
\newcommand{\vp}[1]{\mathbf{p}_{#1}}  % 3D position of VP i
\newcommand{\ed}[1]{\mathbf{d}_{#1}}  % edge direction d_i = V[v1_i] - V[v2_i]

% -------------------------------------------------------------------------
% Steiner / triple-point
% -------------------------------------------------------------------------
\newcommand{\Sv}{\mathbf{S}}           % Steiner (Fermat-Torricelli) point
\newcommand{\Ki}{K_{i}}               % projection matrix (I-n_i n_i^T)/r_i
\newcommand{\Mmat}{M}                 % sum of K_i matrices at a triple point

% -------------------------------------------------------------------------
% Normals and unit vectors
% -------------------------------------------------------------------------
\newcommand{\nhat}{\hat{\mathbf{n}}}  % unit normal to a triangle
\newcommand{\Dhat}{\hat{\boldsymbol{\Delta}}}  % unit segment direction

% -------------------------------------------------------------------------
% Operations
% -------------------------------------------------------------------------
\newcommand{\norm}[1]{\left\|#1\right\|}
\newcommand{\abs}[1]{\left|#1\right|}
\newcommand{\ip}[2]{\langle #1,\, #2 \rangle}

% Projection perpendicular to a unit vector \uhat
\newcommand{\Proj}[1]{\bigl(\mathbf{I} - #1 #1^\top\bigr)}

% argmax operator (winner-take-all assignment in Phase 1 documents)
\DeclareMathOperator*{\argmax}{arg\,max}
\DeclareMathOperator*{\argmin}{arg\,min}

% -------------------------------------------------------------------------
% Calculus shorthands
% -------------------------------------------------------------------------
\newcommand{\pd}[2]{\dfrac{\partial #1}{\partial #2}}
\newcommand{\pdd}[3]{\dfrac{\partial^2 #1}{\partial #2\,\partial #3}}
\newcommand{\grad}{\nabla}
\newcommand{\Hess}[1]{\nabla^2 #1}

% -------------------------------------------------------------------------
% Optimization problem objects
% -------------------------------------------------------------------------
\newcommand{\Preg}{P_{\mathrm{reg}}}  % regular perimeter
\newcommand{\Pst}{P_{\mathrm{st}}}   % Steiner perimeter
\newcommand{\Areg}{A^{\mathrm{reg}}} % regular cell area
\newcommand{\Ast}{A^{\mathrm{st}}}   % Steiner area contribution
\newcommand{\Atgt}{\bar{A}}           % target area per cell
\newcommand{\Lagr}{\mathcal{L}}       % Lagrangian
\newcommand{\objfac}{\sigma}          % IPOPT objective scaling factor

% -------------------------------------------------------------------------
% Code cross-references (for "In the code..." callouts)
% -------------------------------------------------------------------------
\newcommand{\code}[1]{\texttt{#1}}
\newcommand{\codefile}[1]{\texttt{\small #1}}

% -------------------------------------------------------------------------
% Threshold dynamics / approach B (10-mbo-auction-dynamics)
% -------------------------------------------------------------------------
\newcommand{\Dlump}{D}                 % lumped mass matrix diag(v)
\newcommand{\Atau}{A_\tau}             % discrete diffusion operator (M + tau K)^{-1} M
\newcommand{\Etau}{E_\tau}             % heat-content (threshold-dynamics) energy
\newcommand{\Rcell}{R_{\mathrm{cell}}} % radius of the equal-area disc sqrt(Abar/pi)
\newcommand{\hmean}{h_{\mathrm{mean}}} % mean edge length
\newcommand{\hmax}{h_{\mathrm{max}}}   % max edge length
\newcommand{\Gt}{G_t}                  % heat kernel at time t
\newcommand{\Ktau}{K^{\mathrm{res}}_\tau} % resolvent (backward-Euler) kernel
      % loads ALL packages and notation — do not add
                               % packages directly in main.tex

\hypersetup{
  colorlinks = true,
  linkcolor  = blue!60!black,
  citecolor  = green!50!black,
  urlcolor   = blue!60!black,
  pdftitle   = {Phase 1 Energy: Discretization of the Gamma-Convergence Functional},
  pdfauthor  = {surface-partition project}
}

\title{\textbf{Phase 1 Energy: Discretization of the \\ $\Gamma$-Convergence Functional}\\[0.5em]
  \large Dirichlet term, double well, Modica--Mortola limit, and the crispness penalty}
\author{}
\date{July 2026}

\begin{document}
\maketitle

\begin{abstract}
This document derives the discretized Phase~1 relaxation energy minimized by
\code{ProjectedGradientOptimizer} (\codefile{src/optimization/pgd\_optimizer.py}) and
its gradient, exactly as implemented after the double-well correction of
2026-07-08. The continuous functional is the Modica--Mortola-type
$\Gamma$-convergence energy of Bogosel \& Oudet: a Dirichlet interface term plus a
double well, whose $\varepsilon\to0$ limit is the partition perimeter, augmented by a
soft \emph{crispness} penalty. Three quantities are derived from the $P_1$ finite-element
discretization: the Dirichlet term $\varepsilon\,u^\top K u$, the double-well term
$\varepsilon^{-1} q^\top M q$ with $q=u(1-u)$, and the variance penalty
$\lambda\,(1-\mathrm{Var}/T)$; all three gradients are \emph{analytical} and match the
code to machine precision. A remark records the one-token discretization bug this
document's corrected form replaces (the well was previously coded as
$q=u^2(1-u)^2$; see \codefile{docs/reference/phase1\_energy\_discretization\_bug.md}).
\end{abstract}

\tableofcontents
\newpage

\section{Overview}
\label{sec:overview}

Phase~1 of the method (relaxation) replaces the geometric problem ``partition the
surface $S$ into $N$ equal-area regions of minimal total interface length'' by a
smooth minimization over \emph{density fields}. Each region $k$ is represented by a
function $u_k : S \to [0,1]$, where $u_k(x)$ is ``how strongly cell $k$ claims the
point $x$.'' The optimizer minimizes a relaxed energy $E$ over the stacked field
$u = (u_1,\dots,u_N)$ subject to a partition-of-unity constraint
$\sum_k u_k \equiv 1$ and an equal-area constraint $\int_S u_k \,dA = |S|/N$ for
each $k$ (both enforced by projection; \codefile{src/optimization/projection.py}).
As the interface width $\varepsilon \to 0$ the relaxed energy $\Gamma$-converges to
the partition perimeter, so its minimizers approximate minimal-perimeter partitions
\cite{bogosel2023partitions,modica1987gradient}.

This document derives, from the $P_1$ finite-element discretization, the three
additive pieces of $E$ and their gradients \emph{as implemented}. It is the formal
companion to the reference note
\codefile{docs/reference/phase1\_energy\_discretization\_bug.md}, which records that
the double-well term was for a long time mis-discretized (a typo copied from the
source paper) and was corrected on 2026-07-08; the form derived here is the
corrected one. Per the \codefile{docs/math} scope policy, every formula below is a
faithful transcription of the live code.

\section{Notation and setup}
\label{sec:notation}

\begin{remark}[The symbol $\lambda$]
In this document $\lambda$ denotes the \emph{crispness-penalty weight}
(\code{lambda\_penalty}, Section~\ref{sec:penalty}). It is unrelated to the Phase~2
variable-point parameter $\lam$ of the same name used elsewhere in
\codefile{docs/math}; Phase~1 has no variable points.
\end{remark}

\paragraph{Mesh and $P_1$ finite elements.}
$S$ is triangulated with vertices $\V = \{x_1,\dots,x_{|V|}\}$ and triangles $\F$.
Let $\{\varphi_j\}_{j=1}^{|V|}$ be the piecewise-linear ($P_1$) nodal ``hat'' basis,
$\varphi_j(x_i)=\delta_{ij}$. A nodal vector $w\in\mathbb{R}^{|V|}$ defines the $P_1$
function $w_h = \sum_j w_j \varphi_j$, i.e.\ the linear interpolant of the values
$w_j$. The two standard $P_1$ assembly matrices are the \emph{mass} and
\emph{stiffness} matrices
\begin{equation}
  M_{jk} \;=\; \int_S \varphi_j \varphi_k \,dA,
  \qquad
  K_{jk} \;=\; \int_S \grad_\tau \varphi_j \cdot \grad_\tau \varphi_k \,dA,
  \label{eq:MK-def}
\end{equation}
where $\grad_\tau$ is the tangential (surface) gradient. Both are built by
\code{TriMesh.compute\_matrices} (\codefile{src/mesh/tri\_mesh.py}); $M$ and $K$ are
symmetric and sparse. The \emph{lumped mass} vector is the row sum of $M$,
\begin{equation}
  v_j \;=\; \sum_k M_{jk} \;=\; \int_S \varphi_j \, dA \;>\; 0,
  \qquad
  W \;=\; \sum_j v_j \;=\; \int_S 1 \, dA \;=\; |S|,
  \label{eq:lumped-mass}
\end{equation}
so $v_j$ is the area weight of vertex $j$ and $W$ is the total surface area
(\code{self.v}, \code{self.W} in the optimizer). The equal-area target per cell is
$\bar A = W/N$ and the target density mean is $\mu_t = 1/N$.

\paragraph{Interface width.}
The width is tied to the mesh, $\varepsilon = \sqrt{\text{mean triangle area}}
\approx h$ (\codefile{src/pipeline/relaxation.py}, \code{\_setup\_level}), and is
held fixed within a refinement level.

\paragraph{The two exact $P_1$ identities.}
For any nodal vector $w$,
\begin{align}
  \int_S \abs{\grad_\tau w_h}^2 \, dA &\;=\; w^\top K w, \label{eq:dirichlet-identity}\\
  \int_S w_h^2 \, dA &\;=\; w^\top M w, \label{eq:mass-identity}
\end{align}
both immediate from \eqref{eq:MK-def} and bilinearity, since
$w_h = \sum_j w_j \varphi_j$. These are the only two discretization identities the
energy uses; everything else is exact algebra.

\paragraph{The continuous functional.}
Per cell, the relaxed energy is the Modica--Mortola-type functional
\begin{equation}
  F_\varepsilon(u_k)
  \;=\;
  \varepsilon \int_S \abs{\grad_\tau u_k}^2 \, dA
  \;+\;
  \frac{1}{\varepsilon} \int_S W\!\big(u_k\big)\, dA,
  \qquad
  W(s) = s^2 (1-s)^2,
  \label{eq:continuous-Feps}
\end{equation}
with the standard quartic double well $W$ (minima at $s=0$ and $s=1$). The total
Phase~1 energy sums this over the $N$ cells and adds the penalty of
Section~\ref{sec:penalty}:
\begin{equation}
  E(u) \;=\; \sum_{k=1}^{N} F_\varepsilon(u_k) \;+\; \lambda \sum_{k=1}^{N} P(u_k).
  \label{eq:total-continuous}
\end{equation}

\section{The Dirichlet (interface-width) term}
\label{sec:dirichlet}

Discretizing $u_k$ by its $P_1$ interpolant and applying
\eqref{eq:dirichlet-identity} gives the interface-width term exactly:
\begin{equation}
  \varepsilon \int_S \abs{\grad_\tau u_k}^2 \, dA
  \;=\;
  \varepsilon \, u_k^\top K u_k
  \;=\!:\; E^{\mathrm{dir}}_k .
  \label{eq:dirichlet-term}
\end{equation}
Because $K$ is symmetric, the gradient in the nodal unknowns $u_k$ is
\begin{equation}
  \grad_{u_k} E^{\mathrm{dir}}_k \;=\; 2\,\varepsilon \, K u_k .
  \label{eq:dirichlet-grad}
\end{equation}

\begin{tcolorbox}[colback=blue!4!white, colframe=blue!40!black,
                  title={Code: Dirichlet term and gradient}, fontupper=\small]
\codefile{src/optimization/pgd\_optimizer.py}:
\code{compute\_energy} / \code{compute\_gradient} (per-cell loop)\\[2pt]
\code{grad\_term = epsilon * (u.T @ (K @ u))}\\
\code{grad\_grad = 2 * epsilon * (K @ u)}
\end{tcolorbox}

\section{The double-well term (the corrected discretization)}
\label{sec:doublewell}

The interface (binarization) term is
$\varepsilon^{-1}\int_S W(u_k)\,dA = \varepsilon^{-1}\int_S u_k^2(1-u_k)^2\,dA$.
Write the integrand as a perfect square,
\begin{equation}
  W(u_k) \;=\; u_k^2(1-u_k)^2 \;=\; \big[\,u_k(1-u_k)\,\big]^2 \;=\; q_k^2,
  \qquad
  q_k \;:=\; u_k(1-u_k).
  \label{eq:q-def}
\end{equation}
Discretizing the \emph{new} field $q_k$ nodally --- i.e.\ forming the vector with
entries $q_{k,j}=u_{k,j}(1-u_{k,j})$ and taking its $P_1$ interpolant --- and applying
the mass-matrix identity \eqref{eq:mass-identity} gives
\begin{equation}
  \frac{1}{\varepsilon}\int_S W(u_k)\,dA
  \;\approx\;
  \frac{1}{\varepsilon}\int_S (q_k)_h^2 \, dA
  \;=\;
  \frac{1}{\varepsilon}\, q_k^\top M q_k
  \;=\!:\; E^{\mathrm{well}}_k .
  \label{eq:well-term}
\end{equation}

\begin{remark}[Nodal interpolation of the well]
The equality \eqref{eq:mass-identity} is exact, but \eqref{eq:well-term} carries one
approximation: the integrand $q_k=u_k(1-u_k)$ is evaluated at the nodes and then
interpolated linearly, so $(q_k)_h$ is the $P_1$ interpolant of the nodal products,
not the exact quadratic $u_h(1-u_h)$. This is the usual nodal-quadrature treatment of
a polynomial nonlinearity; it is consistent (error $O(h^2)$) and is what the code
computes. Equivalently, the \emph{discrete} double-well energy is defined as
$\varepsilon^{-1}\lVert (q_k)_h \rVert_{L^2}^2$.
\end{remark}

\paragraph{Gradient.}
Since $q_k = u_k \odot (1-u_k)$ acts componentwise, its Jacobian with respect to
$u_k$ is diagonal:
\begin{equation}
  \frac{\partial q_{k,j}}{\partial u_{k,j}}
  \;=\; 1 - 2u_{k,j}
  \qquad\Longrightarrow\qquad
  \frac{\partial q_k}{\partial u_k} \;=\; \operatorname{diag}(1 - 2u_k).
  \label{eq:dq-du}
\end{equation}
With $M$ symmetric, the chain rule applied to $q_k^\top M q_k$ gives
\begin{equation}
  \grad_{u_k} E^{\mathrm{well}}_k
  \;=\;
  \frac{1}{\varepsilon}\,\operatorname{diag}(1-2u_k)\,\big(2 M q_k\big)
  \;=\;
  \frac{2}{\varepsilon}\,\big(M q_k\big)\odot(1-2u_k),
  \label{eq:well-grad}
\end{equation}
where $\odot$ is the Hadamard (elementwise) product.

\begin{tcolorbox}[colback=blue!4!white, colframe=blue!40!black,
                  title={Code: double-well term and gradient (corrected)}, fontupper=\small]
\codefile{src/optimization/pgd\_optimizer.py}:
\code{compute\_energy} / \code{compute\_gradient} (per-cell loop)\\[2pt]
\code{interface\_vec = u * (1 - u)}\hfill\textrm{\small ($=q_k$; not \code{u**2*(1-u)**2})}\\
\code{interface\_term = (1/epsilon) * (interface\_vec.T @ (M @ interface\_vec))}\\
\code{grad\_interface = (2/epsilon) * (M @ interface\_vec) * (1 - 2*u)}
\end{tcolorbox}

\begin{remark}[The discretization bug this form corrects]
\label{rem:bug}
Earlier code set \code{interface\_vec = u**2 * (1-u)**2}, i.e.\ $q_k \mapsto q_k^2$.
That transcribes a \emph{typo} in the source paper \cite{bogosel2023partitions}
(\S3), whose printed definition of the discretized field has an extra squaring, while
its prose target integral $\int u^2(1-u)^2$ and its printed gradient
$2Mv\odot(1-2u)$ both require $v=u(1-u)$. The buggy energy computed
$\varepsilon^{-1}\!\int (u^2(1-u)^2)_h^2 = \varepsilon^{-1}\!\int u^4(1-u)^4$ --- an
eighth-degree well, $\sim\!25\times$ too small in magnitude --- \emph{and} its
gradient \eqref{eq:well-grad} was then no longer the gradient of the coded energy
(finite-difference relative error $\approx 1.2$). With the correct
$q_k=u_k(1-u_k)$, energy \eqref{eq:well-term} and gradient \eqref{eq:well-grad} are
mutually consistent to machine precision (FD relative error $\approx 10^{-8}$). See
\codefile{docs/reference/phase1\_energy\_discretization\_bug.md} for the full audit
and the $\sim\!25\times$ scale shift, which required re-tuning $\lambda$.
\end{remark}

\section{The Modica--Mortola limit: why this is a perimeter}
\label{sec:gamma}

The pairing in \eqref{eq:continuous-Feps} is the classical Modica--Mortola
functional. For a well $W\ge0$ vanishing exactly at $0$ and $1$, the family
$F_\varepsilon$ $\Gamma$-converges as $\varepsilon\to0$ (in $L^1$) to
\begin{equation}
  F_0(u) \;=\; c_W \cdot \operatorname{Per}\big(\{u = 1\}\big),
  \qquad
  c_W \;=\; 2\int_0^1 \sqrt{W(s)}\; ds,
  \label{eq:MM-limit}
\end{equation}
finite only on fields $u\in\{0,1\}$ of finite perimeter \cite{modica1987gradient}. For
the quartic well $W(s)=s^2(1-s)^2$ we have $\sqrt{W(s)} = s(1-s)$ on $[0,1]$, so
\begin{equation}
  c_W \;=\; 2\int_0^1 s(1-s)\,ds
        \;=\; 2\Big[\tfrac12 - \tfrac13\Big]
        \;=\; \tfrac13 .
  \label{eq:cW-value}
\end{equation}
Summing over the $N$ cells as in \eqref{eq:total-continuous}, each interface between
two adjacent regions lies in the boundary of \emph{two} phases and is therefore
counted twice, so the limit of $\sum_k F_\varepsilon(u_k)$ is $2c_W = \tfrac23$ times
the total length of the partition's interface network. The constant is immaterial to
the minimizer: minimizing the relaxed energy at small $\varepsilon$ drives the
partition toward minimal total interface length --- the perimeter problem the method
targets. This is the sense in which the coefficient balance $\varepsilon\,u^\top K u$
versus $\varepsilon^{-1} q^\top M q$ matters: the double-well correction of
Section~\ref{sec:doublewell} restores $W$ to the quartic \eqref{eq:continuous-Feps}
whose limit constant is \eqref{eq:cW-value}, rather than an eighth-degree well with a
different (and much flatter) profile.

\section{The crispness penalty}
\label{sec:penalty}

The double well rewards \emph{pointwise} binarization but does not by itself force a
cell to concentrate its (area-constrained) mass into a compact, sharp region. The
optimizer adds a soft penalty $P$ that rewards each cell's density having the
\emph{variance of a sharp $0/1$ indicator} of the target size. All statistics are
\emph{area-weighted} by the lumped mass $v$ of \eqref{eq:lumped-mass}.

\paragraph{Weighted mean and variance.}
For a cell density $u$ (dropping the cell index),
\begin{equation}
  \mu(u) \;=\; \frac{v^\top u}{W},
  \qquad
  \operatorname{Var}(u) \;=\; \frac{1}{W}\sum_j v_j\,(u_j - \mu)^2
  \;=\; \frac{\big((u-\mu\mathbf 1)\odot v\big)^\top (u-\mu\mathbf 1)}{W}.
  \label{eq:weighted-varmean}
\end{equation}
On the constraint set $v^\top u = \bar A$ the mean is pinned at $\mu = \bar A/W = 1/N
= \mu_t$.

\paragraph{Target variance.}
A sharp indicator equal to $1$ on an area fraction $\mu_t$ and $0$ elsewhere has
area-weighted variance $\mu_t(1-\mu_t)$. The penalty compares $\operatorname{Var}(u)$
to a target
\begin{equation}
  T \;=\;
  \begin{dcases}
    \mu_t(1-\mu_t), & \text{fixed target (\code{penalty\_target\_mode = 'fixed'})},\\
    \mu(1-\mu),     & \text{adaptive target (\code{'adaptive'})},
  \end{dcases}
  \qquad
  T_{\mathrm{eff}} = T + \varepsilon_P,
  \label{eq:target-var}
\end{equation}
with a small floor $\varepsilon_P = $ \code{penalty\_eps} guarding the division. The
penalty per cell is
\begin{equation}
  P(u) \;=\; \lambda\left(1 - \frac{\operatorname{Var}(u)}{T_{\mathrm{eff}}}\right).
  \label{eq:penalty}
\end{equation}
It is minimized (drops toward $\lambda\cdot 0$) as $\operatorname{Var}(u)\to T$, i.e.\
as the cell becomes as ``crisp'' as a sharp indicator of the right size; a diffuse,
low-variance density is penalized. In the code the per-cell penalties are summed and
added to $E$; the weight is $\lambda=$ \code{lambda\_penalty}.

\paragraph{Gradient.}
The gradient of the weighted variance \eqref{eq:weighted-varmean} is
\begin{equation}
  \grad_u \operatorname{Var}(u) \;=\; \frac{2}{W}\,v\odot(u - \mu\mathbf 1),
  \label{eq:grad-var}
\end{equation}
where the terms from the $\mu$-dependence cancel: $\partial\mu/\partial u_k = v_k/W$
and $\sum_j v_j (u_j-\mu)=0$ on any $u$. For the \emph{fixed} target $T$ is constant,
so
\begin{equation}
  \grad_u P \;=\; -\,\frac{\lambda}{T_{\mathrm{eff}}}\,\grad_u \operatorname{Var}(u).
  \label{eq:grad-penalty-fixed}
\end{equation}
For the \emph{adaptive} target $T=\mu(1-\mu)$ depends on $u$ through
$\partial T/\partial u = (1-2\mu)\,v/W$, giving the quotient-rule form
\begin{equation}
  \grad_u P \;=\; -\lambda\!\left[
    \frac{\grad_u \operatorname{Var}(u)}{T_{\mathrm{eff}}}
    \;-\;
    \frac{\operatorname{Var}(u)}{T_{\mathrm{eff}}^{\,2}}\,(1-2\mu)\,\frac{v}{W}
  \right].
  \label{eq:grad-penalty-adaptive}
\end{equation}

\begin{tcolorbox}[colback=blue!4!white, colframe=blue!40!black,
                  title={Code: crispness penalty and gradient}, fontupper=\small]
\codefile{src/optimization/pgd\_optimizer.py}:
\code{compute\_energy} / \code{compute\_gradient} (penalty loop)\\[2pt]
\code{mu = (v @ u) / W};\quad
\code{var\_w = ((center * v) @ center) / W}\hfill\textrm{\small(\code{center = u - mu})}\\
\code{grad\_var = (2/W) * (v * center)}\\
fixed:\ \ \code{G += -lambda * (grad\_var / T\_eff)}\\
adaptive:\ \ \code{G += -lambda * (grad\_var/T\_eff - (var\_w/T\_eff**2)*(1-2*mu)*(v/W))}
\end{tcolorbox}

\begin{remark}[Operating regime after the energy fix]
The $\sim\!25\times$ scale shift from the double-well correction
(Remark~\ref{rem:bug}) rebalances $\lambda$ against the well and Dirichlet terms.
Empirically, on the corrected energy $\lambda$ becomes an effective lever where it was
inert before: a moderate \code{lambda\_penalty}$\approx 5$ resolves the high-$N$
``runt'' cell that the buggy energy manufactured, and seeded initialization becomes
mandatory (random initialization now stalls at the symmetric state). Both effects are
measured in \codefile{docs/experiments/02-corrected-energy-highn-validation/}; they are
consequences of the energy \emph{scale}, not of the formulas derived here.
\end{remark}

\section{Assembled energy and gradient}
\label{sec:assembly}

Collecting Sections~\ref{sec:dirichlet}--\ref{sec:penalty}, with
$q_k = u_k \odot (1-u_k)$, the discretized Phase~1 energy and its gradient (column
$k$ acting on cell $k$) are
\begin{align}
  E(u)
  &\;=\;
  \sum_{k=1}^{N}\Big[\,
    \varepsilon\, u_k^\top K u_k
    \;+\;
    \tfrac{1}{\varepsilon}\, q_k^\top M q_k
  \,\Big]
  \;+\;
  \lambda \sum_{k=1}^{N}\Big(1 - \tfrac{\operatorname{Var}(u_k)}{T_{\mathrm{eff}}}\Big),
  \label{eq:E-assembled}\\[4pt]
  \grad_{u_k} E
  &\;=\;
  \underbrace{2\varepsilon\, K u_k}_{\text{Dirichlet}}
  \;+\;
  \underbrace{\tfrac{2}{\varepsilon}\,(M q_k)\odot(1-2u_k)}_{\text{double well}}
  \;+\;
  \underbrace{\grad_{u_k} P(u_k)}_{\eqref{eq:grad-penalty-fixed}\text{ or }\eqref{eq:grad-penalty-adaptive}}.
  \label{eq:grad-assembled}
\end{align}
The gradient \eqref{eq:grad-assembled} is the unconstrained (Euclidean) gradient the
optimizer then projects onto the tangent of the sum-to-one and equal-area constraints
(\codefile{src/optimization/projection.py}); the constraint handling is outside the
scope of this energy derivation.

\begin{tcolorbox}[colback=blue!4!white, colframe=blue!40!black,
                  title={Code: assembled energy and gradient}, fontupper=\small]
\codefile{src/optimization/pgd\_optimizer.py}:
\code{compute\_energy(x, return\_components)} returns
\code{\{total, grad, interface, penalty\}};
\code{compute\_gradient(x)} returns the stacked $\grad E$ reshaped to
$\mathbb{R}^{|V|\times N}$.
\end{tcolorbox}

\section*{Summary table}
\addcontentsline{toc}{section}{Summary table}

\begin{table}[H]
\centering
\small
\begin{tabular}{@{}p{4.4cm}p{4.0cm}p{2.0cm}p{3.4cm}@{}}
\toprule
Quantity & Discretized form & Method & Implementing function \\
\midrule
Dirichlet term \eqref{eq:dirichlet-term} & $\varepsilon\,u_k^\top K u_k$ & analytical & \code{compute\_energy} \\
\ \ its gradient \eqref{eq:dirichlet-grad} & $2\varepsilon\,K u_k$ & analytical & \code{compute\_gradient} \\
Double-well term \eqref{eq:well-term} & $\varepsilon^{-1} q_k^\top M q_k$, $q_k=u_k(1-u_k)$ & analytical & \code{compute\_energy} \\
\ \ its gradient \eqref{eq:well-grad} & $\tfrac{2}{\varepsilon}(M q_k)\odot(1-2u_k)$ & analytical & \code{compute\_gradient} \\
Weighted variance \eqref{eq:weighted-varmean} & $\tfrac1W \sum_j v_j (u_j-\mu)^2$ & analytical & \code{compute\_energy} \\
\ \ its gradient \eqref{eq:grad-var} & $\tfrac{2}{W} v\odot(u-\mu\mathbf 1)$ & analytical & \code{compute\_gradient} \\
Crispness penalty \eqref{eq:penalty} & $\lambda(1-\operatorname{Var}/T_{\mathrm{eff}})$ & analytical & \code{compute\_energy} \\
\ \ its gradient (fixed / adaptive) & \eqref{eq:grad-penalty-fixed} / \eqref{eq:grad-penalty-adaptive} & analytical & \code{compute\_gradient} \\
$\Gamma$-limit constant \eqref{eq:cW-value} & $c_W = 2\int_0^1\!\sqrt{W} = \tfrac13$ & analytical & --- (theory) \\
\bottomrule
\end{tabular}
\end{table}

\noindent All gradients are analytical and consistent with the corresponding energy
to machine precision (finite-difference relative error $\approx 10^{-8}$). The
formulas above transcribe the corrected code; the superseded (buggy) double-well
discretization is recorded in Remark~\ref{rem:bug} and audited in
\codefile{docs/reference/phase1\_energy\_discretization\_bug.md}.

\bibliographystyle{plain}
\bibliography{../shared/references}

\end{document}
